\section{Elementary incidence and concentration estimates}
\label{sec:elementary}

The statements in this appendix do not depend on the construction.
The incidence identity holds over every finite field and for every set
of points, including a set obtained from earlier random choices. The two
probability estimates apply to a specified filtration or to independent
trials; their applications must still supply the relevant stopping rules
and union bounds.

\begin{lemma}[Affine second moment]\label{lem:affine-second-moment}
Let $q$ be a prime power, let $U\subseteq\F_q^2$, and let $\mathcal L$
be the set of all affine lines. Then
\begin{align*}
 \sum_{\ell\in\mathcal L}|U\cap\ell|&=(q+1)|U|,\\
 \sum_{\ell\in\mathcal L}|U\cap\ell|^2&=|U|^2+q|U|,\\
 \sum_{\ell\in\mathcal L}\left(|U\cap\ell|-\frac{|U|}{q}\right)^2
 &=q|U|-\frac{|U|^2}{q}\le\frac{q^3}{4}.
\end{align*}
For each direction $d$, define
\[
 V_d=\frac1q\sum_{\ell\parallel d}
       \left(|U\cap\ell|-\frac{|U|}{q}\right)^2.
\]
Then $\sum_d V_d=|U|-|U|^2/q^2\le |U|\le q^2$.
\end{lemma}

\begin{proof}
Every point lies on $q+1$ affine lines, proving the first identity.
Every ordered pair of distinct points lies on exactly one line; the
$|U|$ equal pairs each lie on $q+1$ lines. Thus the second sum is
$|U|(|U|-1)+(q+1)|U|$. Expanding the square and using
$|\mathcal L|=q(q+1)$ gives the third identity. Its upper bound follows
by maximizing $qx-x^2/q$ for $0\le x\le q^2$, whose maximum is
$q^3/4$ at $x=q^2/2$. Finally, the parallel classes partition
$\mathcal L$, so division by $q$ gives the directional identity.
\end{proof}

The next estimate is the bounded-increment exponential-moment argument
of Azuma~\cite{Azuma1967}, stated for a process with a predictable drift.
We retain a constant that accommodates centering each increment.

\begin{lemma}[Bounded-drift concentration]
\label{lem:bounded-drift-concentration}
Let $(X_t)_{t=0}^T$ be a real process adapted to a filtration
$(\mathcal F_t)_{t=0}^T$, where $T\ge1$ is an integer. Suppose $L>0$
and $|X_{t+1}-X_t|\le L$ almost surely for $0\le t<T$. Put
$d_t=\E[X_{t+1}-X_t\mid\mathcal F_t]$. For every $u>0$ and each
fixed integer $0\le t\le T$,
\[
 \Prb\left(\left|X_t-X_0-\sum_{j<t}d_j\right|\ge u\right)
 \le2\exp\left(-\frac{u^2}{8TL^2}\right).
\]
\end{lemma}

\begin{proof}
Set $Z_j=X_{j+1}-X_j-d_j$. Then
$\E[Z_j\mid\mathcal F_j]=0$ and $|Z_j|\le2L$. For any real
$\theta$, convexity on $[-2L,2L]$ bounds $e^{\theta Z_j}$ by the
linear interpolation of its endpoint values. Taking conditional
expectations gives
\[
 \E[e^{\theta Z_j}\mid\mathcal F_j]
 \le\cosh(2L\theta)\le\exp(2L^2\theta^2).
\]
The last inequality follows, for example, by integrating
$(\log\cosh x)'=\tanh x\le x$ for $x\ge0$ and using evenness.
Successive conditional expectations therefore give
$\E\exp(\theta\sum_{j<t}Z_j)\le\exp(2TL^2\theta^2)$.
For $\theta>0$, Markov's inequality bounds the upper tail at $u$ by
$\exp(-\theta u+2TL^2\theta^2)$. Choosing
$\theta=u/(4TL^2)$ gives $\exp(-u^2/(8TL^2))$.
Apply the same argument to $-Z_j$ and add the two tails. This proves a
bound at each fixed time; no maximum over times is being taken.
\end{proof}

\begin{lemma}[Independent-trial jumps]
\label{lem:independent-trial-jumps}
Let $Y$ be an integrable real function of independent trials
$\xi_1,\ldots,\xi_T$. Suppose changing only trial $i$ changes $Y$
by at most $c_i$, where $c_i\ge0$. The Doob martingale obtained by
revealing these trials in order has its $i$th increment bounded in
absolute value by $c_i$, and its conditional variance bounded by $c_i^2$.
In particular, if $T\ge1$ and $c_i\le L$ for a constant $L>0$, then
for every $u>0$,
\[
 \Prb(|Y-\E Y|\ge u)\le
 2\exp\left(-\frac{u^2}{8TL^2}\right).
\]
\end{lemma}

\begin{proof}
Fix the values of the first $i-1$ trials. For a possible value $z$ of
trial $i$, average $Y$ over the independent future trials to obtain
$f_i(z)$. Coupling the same future trials for two values $z,z'$ gives
$|f_i(z)-f_i(z')|\le c_i$. Independence ensures that this is the
conditional future law for both values. Thus the range of $f_i$ has
length at most $c_i$. Its deviation from its mean over trial $i$ is
bounded in absolute value by $c_i$; this deviation is exactly the
$i$th Doob increment. Its conditional second moment, and hence its
conditional variance, is at most $c_i^2$. The martingale starts at
$\E Y$, ends at $Y$, and has zero conditional drift. Apply
Lemma~\ref{lem:bounded-drift-concentration} at time $T$.
\end{proof}
